% Generated by exam_latex from exam.md (schema 1.0).
% Compile with XeLaTeX or LuaLaTeX:  xelatex <this file>
\documentclass[UTF8,12pt]{ctexart}

% ---------- Packages ----------
\usepackage{amsmath,amssymb,bm}
\usepackage[a4paper,top=2.8cm,bottom=2.5cm,left=2.4cm,right=2.4cm]{geometry}
\usepackage{enumitem}      % questions / sub-questions / statements
\usepackage{fancyhdr}      % running header and footer
\usepackage{tasks}         % multiple-choice options
\usepackage{needspace}     % keep tall questions on one page
\usepackage{booktabs}      % table rules
\usepackage{graphicx}      % real figures

% ---------- Fonts ----------
% xeCJK does not treat circled digits (① ② ...) as CJK by default, so they
% would come from the Latin font and be missing. Use the CJK font instead.
\ifdefined\xeCJKsetcharclass\xeCJKsetcharclass{"2460}{"24FF}{1}\fi

% ---------- Exam metadata ----------
\newcommand{\examtitle}{2024 普通高等学校招生考试（全国甲卷文）}
\newcommand{\examsubtitle}{普通高等学校招生考试（全国甲卷文）}
\newcommand{\examheaderleft}{2024 普通高等学校招生考试}
\newcommand{\examheaderright}{全国甲卷文}

% ---------- Page style ----------
\pagestyle{fancy}
\fancyhf{}
\setlength{\headheight}{24.1pt}
\renewcommand{\headrulewidth}{0pt}
\fancyhead[L]{\small\examheaderleft}
\fancyhead[R]{\small\examheaderright}
\fancyfoot[C]{\small\thepage}

\setlength{\parindent}{0pt}
\setlength{\parskip}{0.5em}
\linespread{1.25}
\renewcommand{\arraystretch}{1.3}   % roomier tables

% ---------- Section headings ----------
\ctexset{section={format=\large\bfseries, beforeskip=1.6em, afterskip=0.8em}}

% ---------- Lists ----------
\newlist{questions}{enumerate}{1}
\setlist[questions]{label=\textbf{\arabic*.},
  leftmargin=2.2em, itemsep=1.1em, topsep=0pt, parsep=0.4em}

\newlist{subquestions}{enumerate}{2}
\setlist[subquestions,1]{label=(\arabic*), leftmargin=2.6em,
  itemsep=0.3em, topsep=0.2em, parsep=0pt}
\setlist[subquestions,2]{label=(\alph*), leftmargin=2.4em,
  itemsep=0.2em, topsep=0.1em, parsep=0pt}

\newlist{statements}{description}{1}   % ① ② ③ ④ propositions
\setlist[statements]{font=\normalfont, align=left, leftmargin=3em,
  labelwidth=1.6em, labelsep=0.4em, itemsep=0.1em, topsep=0.2em, parsep=0pt}

% ---------- Multiple-choice options ----------
\settasks{label=(\Alph*), label-width=1.8em, column-sep=1em,
  before-skip=0.2em, after-skip=0.2em, after-item-skip=0.2em}

% ---------- Macros ----------
\newcommand{\blank}{\underline{\hspace{4em}}}   % fill-in blank
\newcommand{\examfigure}[3]{%   width (cm), height (cm), label
  \parbox[c][#2cm][c]{#1cm}{\centering\small\textit{[ #3 ]}}}
\newcommand{\makeexamtitle}{%
  \begin{center}
    {\LARGE\bfseries\examtitle}\par
    \vspace{0.4em}{\Large\bfseries\examsubtitle}\par
  \end{center}\vspace{0.4em}}

\begin{document}

\makeexamtitle

%==================================================================
\section*{一、单选题}
%==================================================================

\begin{questions}

  % Question 1
  \item 若 $z = \sqrt{2}i$，则 $z\overline{z} =$ （ ）

    \begin{tasks}(4)
      \task $-2$
      \task $-\sqrt{2}$
      \task $\sqrt{2}$
      \task $2$
    \end{tasks}

  % Question 2
  \item 已知集合 $A = \{1, 2, 3, 4, 5, 9\}$, $B = \{x \mid x + 1 \in A\}$, 则
    $A \cap B =$ （ ）

    \begin{tasks}(2)
      \task $\{1, 2, 3\}$
      \task $\{3, 4, 9\}$
      \task $\{1, 2, 3, 4\}$
      \task $\{2, 3, 4, 5\}$
    \end{tasks}

  % Question 3
  \item 设向量 $\boldsymbol{a} = (x + 1, x)$, $\boldsymbol{b} = (x, 2)$,
    $\boldsymbol{c} = (1, 1)$，若 $a - 2$b 与 $c$ 共线，则 $x$ = （ ）

    \begin{tasks}(4)
      \task $\frac{5}{2}$
      \task $\frac{1}{2}$
      \task $-\frac{1}{2}$
      \task $-\frac{5}{2}$
    \end{tasks}

  % Question 4
  \item 某独唱比赛的决赛阶段共有甲、乙、丙、丁四人参加，每人出场一次，
    出场次序由随机抽签确定。则丙不是第一个出场，且甲或乙最后出场的概率是（ ）

    \begin{tasks}(4)
      \task $\frac{1}{6}$
      \task $\frac{1}{4}$
      \task $\frac{1}{3}$
      \task $\frac{1}{2}$
    \end{tasks}

  % Question 5
  \item 记 $S_{n}$ 为等差数列 $\{a_{n}\}$ 的前 $n$ 项和. 已知 $S_{9}=1$, 则
    $a_{3}+a_{7}=$ （ ）

    \begin{tasks}(4)
      \task $\frac{1}{9}$
      \task $\frac{2}{9}$
      \task $\frac{1}{3}$
      \task $\frac{2}{3}$
    \end{tasks}

  % Question 6
  \item 已知双曲线的两个焦点分别为 $(0,4)$， $(0,-4)$，点 $(-6,4)$ 在该双曲线上，
    则该双曲线的离心率为（ ）

    \begin{tasks}(4)
      \task $4$
      \task $3$
      \task $2$
      \task $\sqrt{2}$
    \end{tasks}

  % Question 7
  \item 设函数 $f(x)=\frac{e^{x}+2\sin x}{1+x^{2}}$ ，则曲线 $y=f(x)$ 在点 $(0,1)$
    处的切线与两坐标轴所围成的三角形的面积为 （ ）

    \begin{tasks}(4)
      \task $\frac{1}{6}$
      \task $\frac{1}{3}$
      \task $\frac{1}{2}$
      \task $\frac{2}{3}$
    \end{tasks}

  % Question 8
  \Needspace{9cm}
  \item 函数 $y = -x^{2} + \left(e^{x} - e^{-x}\right)\sin x$ 在区间 $[-2.8, 2.8]$
    的图象大致为（ ）

    \begin{tasks}(2)
      \task \IfFileExists{imgs/img_in_image_box_177_1121_441_1323.jpg}{\includegraphics[width=3.3cm]{imgs/img_in_image_box_177_1121_441_1323.jpg}}{\examfigure{3.3}{2.52}{Option A}}
      \task \IfFileExists{imgs/img_in_image_box_512_1126_779_1332.jpg}{\includegraphics[width=3.3cm]{imgs/img_in_image_box_512_1126_779_1332.jpg}}{\examfigure{3.3}{2.55}{Option B}}
      \task \IfFileExists{imgs/img_in_image_box_177_1370_440_1556.jpg}{\includegraphics[width=3.3cm]{imgs/img_in_image_box_177_1370_440_1556.jpg}}{\examfigure{3.3}{2.33}{Option C}}
      \task \IfFileExists{imgs/img_in_image_box_516_1353_785_1566.jpg}{\includegraphics[width=3.3cm]{imgs/img_in_image_box_516_1353_785_1566.jpg}}{\examfigure{3.3}{2.61}{Option D}}
    \end{tasks}

  % Question 9
  \item 直线 $2x - y - 2 = 0$ 与圆 $x^{2} + y^{2} - 6x - 8y = 0$ 交于 $A, B$ 两点，则
    $|AB|$ = （ ）

    \begin{tasks}(4)
      \task $4$
      \task $5$
      \task $8$
      \task $10$
    \end{tasks}

  % Question 10
  \item 已知 $\frac{\cos\alpha}{\cos\alpha-\sin\alpha}=\sqrt{3}$，则
    $\tan\left(\alpha+\frac{\pi}{4}\right)=$ （ ）

    \begin{tasks}(4)
      \task $2\sqrt{3}+1$
      \task $2\sqrt{3}-1$
      \task $\frac{\sqrt{3}}{2}$
      \task $1-\sqrt{3}$
    \end{tasks}

  % Question 11
  \item 设 $\alpha, \beta$ 为两个平面, $m, n$ 为两条直线, 且 $\alpha \cap \beta = m$.
    下述四个命题:

    \begin{statements}
      \item[①] 若 $m \parallel n$, 则 $n \parallel \alpha$ 或 $n \parallel \beta$;
      \item[②] 若 $m \perp n$, 则 $n \perp \alpha$ 或 $n \perp \beta$;
      \item[③] 若 $n \parallel \alpha$ 且 $n \parallel \beta$, 则 $m \parallel n$;
      \item[④] 若 $n$ 与 $\alpha, \beta$ 所成的角相等, 则 $m \perp n$.
    \end{statements}

  其中所有真命题的编号是（ ）

    \begin{tasks}(4)
      \task ①③
      \task ②④
      \task ①②③
      \task ①③④
    \end{tasks}

  % Question 12
  \item 记 $\triangle ABC$ 的内角 $A, B, C$ 的对边分别为 $a, b, c$, 已知
    $B = 60^{\circ}$, $b^{2} = \frac{9}{4}ac$, 则 $\sin A + \sin C =$ （ ）

    \begin{tasks}(4)
      \task $\frac{3}{2}$
      \task $\sqrt{2}$
      \task $\frac{\sqrt{7}}{2}$
      \task $\frac{\sqrt{3}}{2}$
    \end{tasks}

\end{questions}

%==================================================================
\section*{二、填空题}
%==================================================================

\begin{questions}[start=13]

  % Question 13
  \item 函数 $f(x) = \sin x - \sqrt{3} \cos x$ 在区间 $[0, \pi]$ 的最大值为 \blank。

  % Question 14
  \item 已知圆台甲、乙的上底面半径均为 $r_1$，下底面半径均为 $r_2$，圆台甲、
    乙的母线长分别为 $2(r_2 - r_1)$， $3(r_2 - r_1)$，则圆台甲与乙的体积之比为\blank。

  % Question 15
  \item 已知 $a > 1$ 且 $\frac{1}{\log_8 a} - \frac{1}{\log_a 4} = -\frac{5}{2}$，则
    $a =$\blank。

  % Question 16
  \item 当 $x > 0$ 时，曲线 $y = x^{3} - 3x$ 与曲线 $y = -(x - 1)^{2} + a$ 有两个交点，
    则 $a$ 的取值范围是 \blank。

\end{questions}

%==================================================================
\section*{三、解答题}
%==================================================================

\begin{questions}[start=17]

  % Question 17
  \item 记 $S_{n}$ 为等比数列 $\{a_{n}\}$ 的前 $n$ 项和，已知 $2S_{n}=3a_{n+1}-3$.

    \begin{subquestions}
      \item
      求 $\{a_{n}\}$ 的通项公式;
      \item
      求数列 $\{S_{n}\}$ 的前 $n$ 项和.
    \end{subquestions}

  % Question 18
  \item 某工厂进行生产线智能化升级改造. 升级改造后, 从该工厂甲、
    乙两个车间的产品中随机抽取 150 件进行检验, 数据如下:

    \begin{center}
      \begin{tabular}{lcccc}
        \toprule
               & 优级品 & 合格品 & 不合格品 & 总计 \\
        甲车间 & $26$   & $24$   & $0$      & $50$ \\
        乙车间 & $70$   & $28$   & $2$      & $100$ \\
        总计   & $96$   & $52$   & $2$      & $150$ \\
        \bottomrule
      \end{tabular}
    \end{center}

    \begin{subquestions}
      \item
      填写如下列联表:

        \begin{center}
          \begin{tabular}{lcc}
            \toprule
                   & 优级品 & 非优级品 \\
            甲车间 &        &  \\
            乙车间 &        &  \\
            \bottomrule
          \end{tabular}
        \end{center}

      能否有 95\% 的把握认为甲、乙两车间产品的优级品率存在差异？能否有 99\%
        的把握认为甲、乙两车间产品的优级品率存在差异？
      \item
      已知升级改造前该工厂产品的优级品率 $p = 0.5$. 设 $\overline{p}$ 为升级改造后抽取的
        $n$ 件产品的优级品率，如果 $\overline{p} > p + 1.65\sqrt{\frac{p(1 - p)}{n}}$，
        则认为该工厂产品的优级品率提高了。根据抽取的 150 件产品的数据，
        能否认为生产线智能化升级改造后，该工厂产品的优级品率提高了？
        $(\sqrt{150} \approx 12.247)$
    \end{subquestions}

    附： $K^{2} = \frac{n(ad - bc)^{2}}{(a + b)(c + d)(a + c)(b + d)}$

    $\frac{P(K^{2} \geq k)}{k} | 0.050$ 0.010 0.001

    3.841 6.635 10.828

  % Question 19
  \item 如图，在以 $A, B, C, D, E, F$ 为顶点的五面体中，四边形 $ABCD$ 与四边形 $ADEF$
    均为等腰梯形，%
    $EF \parallel AD, BC \parallel AD, AD = 4, AB = BC = EF = 2, ED = \sqrt{10}, FB = 2 \sqrt{3}, M$
    为 $AD$ 的中点.

    \begin{center}
      \IfFileExists{imgs/img_in_image_box_2059_1248_2291_1465.jpg}{\includegraphics[width=6.2cm]{imgs/img_in_image_box_2059_1248_2291_1465.jpg}}{\examfigure{6.2}{5.8}{Figure}}
    \end{center}

    \begin{subquestions}
      \item
      证明：$BM\parallel$平面$CDE$；
      \item
      求 $M$ 到平面 $FAB$ 的距离.
    \end{subquestions}

  % Question 20
  \item 已知函数 $f(x) = a(x - 1) - \ln x + 1$.

    \begin{subquestions}
      \item
      讨论 $f(x)$ 的单调性;
      \item
      设 $a \leqslant 2$，证明：当 $x > 1$ 时， $f(x) < \mathrm{e}^{x-1}$.
    \end{subquestions}

  % Question 21
  \item 设椭圆 $C: \frac{x^{2}}{a^{2}}+\frac{y^{2}}{b^{2}}=1 (a > b > 0)$ 的右焦点为
    $F$, 点 $M \left(1,\frac{3}{2}\right)$ 在 $C$ 上, 且 $MF \perp x$ 轴.

    \begin{subquestions}
      \item
      求 $C$ 的方程;
      \item
      过点 $P(4,0)$ 的直线交 $C$ 于 $A, B$ 两点，$N$ 为线段 $FP$ 的中点，直线 $NB$
        交直线 $MF$ 于点 $Q$，证明： $AQ \perp y$ 轴.
    \end{subquestions}

  % Question 22
  \item 在直角坐标系 $xOy$ 中，以坐标原点为极点，$x$ 轴正半轴为极轴建立极坐标系，曲线
    $C$ 的极坐标方程为 $\rho = \rho \cos \theta + 1$.

    \begin{subquestions}
      \item
      写出 $C$ 的直角坐标方程;
      \item
      设直线 $l:$ \[\begin{cases}x=t, \\y=t+a\end{cases}\] (t 为参数)，若 $C$ 与 $l$
        相交于 $A, B$ 两点，且 $|AB|=2$，求 $a$.
    \end{subquestions}

  % Question 23
  \item 已知实数 $a, b$ 满足 $a + b \geqslant 3$.

    \begin{subquestions}
      \item
      证明： $2a^{2}+2b^{2}>a+b;$
      \item
      证明： $\left|a-2b^{2}\right|+\left|b-2a^{2}\right|\geqslant6$
    \end{subquestions}

\end{questions}

\end{document}
